% ECNU 模板核心设置（供 Quarto PDF 使用）

\usepackage{graphicx}
\usepackage{fancyhdr}
\usepackage{caption}
\usepackage{bicaption}
\usepackage{titletoc}
\usepackage{enumitem}
\usepackage{fontspec}
\usepackage{amssymb}
\usepackage{etoolbox}

% 仅在目录中保留中文图表条目
\captionsetup[bi-second]{list=off}
% 英文双语标题在不同环境下使用 Figure / Table
\AtBeginEnvironment{figure}{\captionsetup[bi-second]{name=Figure}}
\AtBeginEnvironment{table}{\captionsetup[bi-second]{name=Table}}
\usepackage{geometry}
\usepackage{hyperref}

% 打印模式开关（默认关闭）
% 需要打印时取消下一行注释即可：
% \def\PrintMode{}
\ifdefined\PrintMode
  \def\SideMode{twoside}
  \def\ClearPageStyle{\cleardoublepage}
  \makeatletter\@twosidetrue\makeatother
\else
  \def\SideMode{oneside}
  \def\ClearPageStyle{\clearpage}
  \makeatletter
  \@twosidefalse
  \let\cleardoublepage\clearpage
  \makeatother
\fi

% 版芯控制（参考 ECNU LaTeX 模板）
\setlength{\oddsidemargin}{0.57cm}
\voffset-6mm
\textwidth=150mm
\textheight=230mm
\headwidth=150mm

% 打印模式下镜像版心（偶数页左右边距对调）
\ifdefined\PrintMode
  \setlength{\evensidemargin}{\dimexpr\paperwidth-\textwidth-2in-\oddsidemargin\relax}
\else
  \setlength{\evensidemargin}{\oddsidemargin}
\fi

% 段落缩进与段前后距（避免 parskip 拉大间距）
\setlength{\parindent}{2em}
\setlength{\parskip}{0pt}

% 模板字号宏（封面与部分页会用到）
\newcommand{\yihao}{\fontsize{26pt}{36pt}\selectfont}
\newcommand{\erhao}{\fontsize{22pt}{28pt}\selectfont}
\newcommand{\xiaoer}{\fontsize{18pt}{18pt}\selectfont}
\newcommand{\sanhao}{\fontsize{16pt}{24pt}\selectfont}
\newcommand{\xiaosan}{\fontsize{15pt}{22pt}\selectfont}
\newcommand{\sihao}{\fontsize{14pt}{21pt}\selectfont}
\newcommand{\banxiaosi}{\fontsize{13pt}{19.5pt}\selectfont}
\newcommand{\xiaosi}{\fontsize{12pt}{18pt}\selectfont}
\newcommand{\dawuhao}{\fontsize{11pt}{11pt}\selectfont}
\newcommand{\wuhao}{\fontsize{10.5pt}{15.75pt}\selectfont}

% 列表间距（与模板保持接近）
\setenumerate{itemsep=0pt,partopsep=0pt,parsep=\parskip,topsep=0pt}
\setitemize{itemsep=0pt,partopsep=0pt,parsep=\parskip,topsep=0pt}

% 图表标题分隔符：图 1.1 标题（空格）
\newcommand{\loflabel}{图}
\newcommand{\lotlabel}{表}
\captionsetup[table]{labelsep=quad}
\captionsetup[figure]{labelsep=quad}

% 图表编号按章节，并用 1-1 形式
\makeatletter
\@addtoreset{figure}{chapter}
\@addtoreset{table}{chapter}
\@addtoreset{equation}{chapter}
\makeatother
\renewcommand{\thefigure}{\thechapter-\arabic{figure}}
\renewcommand{\thetable}{\thechapter-\arabic{table}}
% 公式编号按章节，形如 (3-1)；equation 环境本身居中排版，编号自动顶右
\renewcommand{\theequation}{\thechapter-\arabic{equation}}

% 页眉（与模板保持接近）
\pagestyle{fancy}
\fancyhead{}
\setlength{\headheight}{24pt}
\fancyhead[L]{\small\parbox[t]{0.45\textwidth}{\raggedright 华东师范大学20XX届博士学位论文}}
\fancyhead[R]{\small\parbox[t]{0.5\textwidth}{\raggedleft \leftmark}}
\renewcommand{\headrulewidth}{1pt}
\renewcommand{\chaptermark}[1]{\markboth{第\thechapter 章\ #1}{}}

% macOS 下 ctex 的 \heiti 默认映射到「Heiti SC Light」（细体），
% 明显比 Windows 的黑体（SimHei，本身即为粗壮的单一字重）纤细，
% 导致目录/标题的黑体效果不明显。改用「Heiti SC Medium」作为 \heiti 的默认字重。
\IfFontExistsTF{Heiti SC}{\setCJKfamilyfont{zhhei}{Heiti SC Medium}}{}

% 目录样式（使用 ctex 的 \zihao 确保中文字号生效）
% 一级目录：黑体，四号；二级及以下：黑体，小四
% \dottedcontents 最后一个参数是 \titlerule* 里每个点号重复占用的宽度，
% 数值越大点号间距越疏（此前 9.5pt 显得像拉长的虚线），改小以获得紧密的点线效果
\dottedcontents{chapter}[1.5cm]{\addvspace{6pt}\zihao{4}\heiti}{3.8em}{4pt}
\dottedcontents{section}[1.5cm]{\zihao{-4}\heiti}{2.8em}{4pt}
\dottedcontents{subsection}[1.5cm]{\zihao{-4}\heiti}{3.8em}{4pt}

% 图表目录样式：宋体，小四
\dottedcontents{figure}[1.5cm]{\zihao{-4}}{5em}{4pt}
\dottedcontents{table}[1.5cm]{\zihao{-4}}{5em}{4pt}

% 去除图表目录中不同章节间的额外间距
\CTEXsetup[lofskip=0pt, lotskip=0pt]{chapter}

% 章节标题前后间距（参考 ECNU LaTeX 模板）
\CTEXsetup[beforeskip = 0pt]{chapter}
\CTEXsetup[afterskip = 20pt]{chapter}

% 英文字体：优先 Times New Roman，不存在则回退
\IfFontExistsTF{Times New Roman}{\setmainfont{Times New Roman}}{\setmainfont{TeX Gyre Termes}}

% 链接颜色与引文样式在文档开始时统一覆盖
\AtBeginDocument{
  \makeatletter
  \expandafter\def\csname bibstyle@GBT7714-2005\endcsname{\NAT@numbers}
  \expandafter\def\csname bibstyle@GBT7714-2005NLang\endcsname{\NAT@numbers}
  \makeatother
  \hypersetup{
    colorlinks=true,
    linkcolor=black,
    urlcolor=black,
    citecolor=black,
    bookmarksnumbered=true
  }
  \setcitestyle{authoryear}
  \renewcommand{\bibsection}{\chapter*{\bibname}}
  \setlength{\bibhang}{2em}
}

% 让 plain 样式继承 fancy 页眉
\makeatletter
\let\ps@plain\ps@fancy
\makeatother

% 禁用 Quarto 默认标题页，避免封面前多出一页
\renewcommand{\maketitle}{}

% 封面/声明页宏（请按需修改）
\def\thesisTitle{这里定义标题}
\def\thesisTitleNoWrap{这里还是标题，不换行的标题，用于原创性声明}
\def\thesisETitle{Title}
